% Part: lambda-calculus
% Chapter: syntax
% Section: terms

\documentclass[../../../include/open-logic-section]{subfiles}

\begin{document}

\olfileid{lam}{syn}{trm}
\olsection{পদ}

ল্যাম্বডা ক্যালকুলাসের পদগুলি চলরাশি $\Obj{v_0}$, $\Obj{v_1}$, \dots-এর
অসীম ভাণ্ডার, ``$\lambd$'' সংকেত এবং বন্ধনী থেকে আরোহীভাবে গড়ে ওঠে।
চলরাশি নির্দেশ করতে আমরা $x$, $y$, $z$, \dots{} এবং পদ নির্দেশ করতে
$M$, $N$, $P$, \dots{} ব্যবহার করব।

\begin{defn}[পদ] \ollabel{defn:term}
ল্যাম্বডা ক্যালকুলাসের \emph{পদসমূহের} সেটটি নিচের নিয়মে আরোহীভাবে
সংজ্ঞায়িত:
\begin{enumerate}
  \item \ollabel{defn:term-var} $x$ চলরাশি হলে $x$ একটি পদ।
  \item \ollabel{defn:term-abs} $x$ চলরাশি এবং $M$ পদ হলে
    $(\lambd[x][M])$ একটি পদ।
  \item \ollabel{defn:term-app} $M$ ও $N$ উভয়ই পদ হলে $(MN)$ একটি পদ।
\end{enumerate}
\end{defn}

কোনো পদ $(\lambd[x][M])$ যদি \olref{defn:term-abs} অনুসারে গঠিত হয়,
তাহলে বলি সেটি একটি \emph{বিমূর্তনের} ফল, এবং $\lambd[x]$-এর~$x$-কে
একটি \emph{!!{parameter}} বলা হয়। \olref{defn:term-app} অনুসারে গঠিত
কোনো পদ $(MN)$ হলো একটি \emph{প্রয়োগের} ফল।

উপরে সংজ্ঞায়িত পদগুলি পূর্ণ-বন্ধনীযুক্ত। এতে লেখা বেশ দুরূহ হয়ে উঠতে
পারে, যেমন $(\lambd[x][((\lambd[x][x])(\lambd[x][(xx)]))])$ পদটি দেখায়।
বন্ধনী এড়ানোর কিছু প্রথা আমরা চালু করব। তবে আনুষ্ঠানিক সংজ্ঞাটি ব্যবহার
করে \olref{defn:term} অনুসারে কোনো পদ কীভাবে গঠিত, তা সহজে নির্ধারণ করা
যায়। যেমন, $(\lambd[x][((\lambd[x][x])(\lambd[x][(xx)]))])$ পদটি
গঠনের শেষ ধাপটি অবশ্যই এমন একটি বিমূর্তন, যার !!{parameter} হলো~$x$।
এই বিমূর্তনটি $((\lambd[x][x])(\lambd[x][(xx)]))$ পদ থেকে পাওয়া যায়;
আর সেই পদটি দুটি পদের প্রয়োগ। ওই পদদুটির প্রত্যেকটিই আবার একটি
বিমূর্তনের ফল, এবং এভাবেই বিশ্লেষণ চলতে থাকে।

\begin{prob}
$(\lambd[g][(\lambd[x][(g (x x))])
  (\lambd[x][(g (x x))])])$-এর গঠন বর্ণনা করো।
\end{prob}

\end{document}
